
利用第~\ref{subsec:5-3-2} 小节所述的评价映射，
得到
\begin{equation}
\begin{aligned}
\operatorname{ev}\colon{}
&\overset{\circ}{\mathcal M}_{k_1,k_0}
\bigl(\{J_{\tau,t}\}_{\tau,t};
[h,w],[h',w'];\operatorname{top}(\tau)\bigr)
\\
&\longrightarrow
\bigl([0,1]^{k_1}\times(L^{(0)})^{k_1}\bigr)
\times
\bigl([0,1]^{k_0}\times(L^{(0)})^{k_0}\bigr).
\end{aligned}
\tag{7.2.373}\label{eq:7-2-373}
\end{equation}
这里暂按源书保留两个 $L^{(0)}$ 因子
\translatornote{源书 (7.2.373)--(7.2.375)
均把第 $k_1$ 组标记点的空间因子
写成 $(L^{(0)})^{k_1}$；
第 5.3.2 小节的同型评价映射及上边界
标记点的型别则指向 $(L^{(1)})^{k_1}$。
本页后续输入链还涉及加撇的 $L^{(i)\prime}$，
故暂保留源式并将整体型别列为待核。}。
注意，在第~\ref{subsec:5-3-2} 小节，
$\operatorname{ev}_{\mathrm{time},(i),j}$
（$i=0,1$）定义为取值于 $\mathbb R$
的映射。将它与 $\chi$ 复合
\translatornote{源书此处写与 $\chi^{-1}$ 复合；
但第 5.3.2 小节定义
$\chi\colon\mathbb R\to(0,1)$，
故从 $\mathbb R$ 值时间评价得到
$[0,1]$ 值坐标须复合 $\chi$，
而非 $\chi^{-1}$。}，
便得到上述评价映射的 $[0,1]$ 因子。
评价映射 $\operatorname{ev}$
可延拓至模空间 (7.2.370)：
若相应标记点位于滑向
$\tau\to-\infty$ 或 $\tau\to+\infty$
的分量上，就分别取其 $[0,1]$ 分量为
$0$ 或 $1$。

注意，$\operatorname{ev}$ 对
\begin{equation}
\begin{aligned}
\partial\Bigl(
&\bigl([0,1]^{k_1}\times(L^{(0)})^{k_1}\bigr)
\times
\bigl([0,1]^{k_0}\times(L^{(0)})^{k_0}\bigr)
\Bigr)
\end{aligned}
\tag{7.2.374}\label{eq:7-2-374}
\end{equation}
的原像，落在模空间 (7.2.370) 的
余维一边界内。
因此，我们给模空间添加一个领状邻域，
并把映射 (7.2.373)
\translatornote{源书这里写“延拓映射 (7.2.173)”，
但 (7.2.173) 是注 7.2.172 中的纤维积，
并非本页的评价映射；此处按紧前的
(7.2.373) 定型。}
延拓为
\begin{equation}
\begin{aligned}
\operatorname{ev}\colon{}
&\mathcal M_{k_1,k_0}^{\operatorname{twp}(\tau)+}
\bigl(\{J_{\tau,t}\}_{\tau,t};
[h,w],[h',w'];\operatorname{top}(\tau)\bigr)
\\
&\longrightarrow
\bigl([0,1]^{k_1}\times(L^{(0)})^{k_1}\bigr)
\times
\bigl([0,1]^{k_0}\times(L^{(0)})^{k_0}\bigr).
\end{aligned}
\tag{7.2.375}\label{eq:7-2-375}
\end{equation}
在 (7.2.375) 中，符号 $+$ 表示
给模空间 (7.2.370) 添加领状邻域。
可以假设 $\operatorname{ev}$
在这一领状邻域上为积型映射。

另有评价映射
\begin{equation}
\begin{aligned}
&(\operatorname{ev}_{-\infty},
\operatorname{ev}_{+\infty})\colon
\mathcal M_{k_1,k_0}^{\operatorname{twp}(\tau)+}
\bigl(\{J_{\tau,t}\}_{\tau,t};
[h,w],[h',w'];\operatorname{top}(\tau)\bigr)
\\
&\hspace{38mm}\longrightarrow R_h\times R'_{h'}.
\end{aligned}
\tag{7.2.376}\label{eq:7-2-376}
\end{equation}
此处正无穷端的靶写为 $R'_{h'}$
\translatornote{源书 (7.2.376) 的第二因子未加撇，
写作 $R_{h'}$；但图式 7.2.38 与
第 5.3.2 小节的正无穷端评价靶均为
$R'_{h'}$，故依端点模型补撇。}。

现取 $\mathcal P_j^{(i)}$（$i=0,1$）
为 $[0,1]_\rho\times L^{(i)\prime}$ 上的链，
并取 $S$ 为 $R_h$ 上的链
\translatornote{源书明确将 $\mathcal P_j^{(i)}$
置于 $[0,1]_\rho\times L^{(i)\prime}$；
但图式 7.2.40 的中间链模型及
(7.2.373)--(7.2.375) 的评价靶写的是未加撇的
$L^{(i)}$（首因子还混写为 $L^{(0)}$）。
此处保留原样，所需端点推前或模型同一
尚待核明。}。
利用 (7.2.375) 的映射和
(7.2.376) 中的
$\operatorname{ev}_{-\infty}$
取纤维积，得到
\begin{equation}
\begin{aligned}
\mathcal M_{k_1,k_0}^{\operatorname{twp}(\tau)+}
\bigl(&\{J_{\tau,t}\}_{\tau,t};
[h,w],[h',w'];\operatorname{top}(\tau);\\
&\vec{\mathcal P}^{(1)},S,
\vec{\mathcal P}^{(0)}\bigr).
\end{aligned}
\tag{7.2.377}\label{eq:7-2-377}
\end{equation}
注意，在 (5.3.24) 中，
我们用 $L^{(i)\prime}$ 上的链
$P_j^{(i)}$ 取纤维积；
而这里用的是
$[0,1]_\rho\times L^{(i)\prime}$
上的链。

取适当扰动 $\mathfrak s$ 后，映射
\[
\begin{aligned}
\vec{\mathcal P}^{(1)}\otimes S
\otimes\vec{\mathcal P}^{(0)}
\longmapsto
(\operatorname{ev}_{+\infty})_*
\Bigl(
\mathcal M_{k_1,k_0}^{\operatorname{twp}(\tau)+}
\bigl(&\{J_{\tau,t}\}_{\tau,t};
[h,w],[h',w'];\operatorname{top}(\tau);\\
&\vec{\mathcal P}^{(1)},S,
\vec{\mathcal P}^{(0)}\bigr)^{\mathfrak s}
\Bigr)
\end{aligned}
\]
定义了 (7.2.368) 的 $A_\infty$ 双模同态。
我们可以用模空间 (7.2.377) 证明
定理~\ref{thm:5-3-14}，
不必使用第~\ref{subsec:5-3-4} 小节
所用的模空间 (5.3.24)。
这就给出定理~\ref{thm:5-3-14} 的另一种证明。
第 5.3 节其他定理的证明也可适当修改，
改用这里的替代同态。
